Las dos series dan una línea recta en su gráfica de senos, y eso por sí solo ya confirma la ley de Snell: el cociente entre los senos se mantiene constante sin importar el ángulo de entrada.

Pero los dos índices no salieron iguales.

Con la luz entrando por la cara curva se obtuvo $n = 1.494$, y comparado con el valor de referencia del acrílico el error es
\begin{equation}
E_{curva} = \frac{\left| 1.49 - 1.494 \right|}{1.49} \times 100 = 0.2 ,
\label{eq:error_curva}
\end{equation}
es decir, 0.2 por ciento, un resultado excelente. Con la luz entrando por la cara plana, en cambio, salió $n = 1.351$, con un error de
\begin{equation}
E_{plana} = \frac{\left| 1.49 - 1.351 \right|}{1.49} \times 100 = 9.3 ,
\label{eq:error_plana}
\end{equation}
que equivale a 9.3 por ciento.

¿Por qué tanta diferencia si la lente es la misma? La respuesta está en la tabla \ref{tab:plana}. Todos los ángulos de refracción de esa serie quedaron algo grandes, y por eso los puntos de la figura \ref{fig:plana} caen siempre debajo de la línea de referencia. Un error así, que empuja todos los datos hacia el mismo lado, no se debe al azar, sino a algo del montaje. Lo más probable es que, al usar esa cara, el centro de la cara plana no quedara exactamente en el centro de la mesa. Si eso pasa, el rayo ya no viaja dentro de la lente a lo largo de un radio, vuelve a doblarse al salir por la cara curva y la lectura en la escala ya no es el ángulo de refracción verdadero. El efecto es más fuerte en los ángulos chicos, donde unos cuantos grados de diferencia pesan mucho; en la tabla, los índices entre 10 y 25 grados rondan 1.24, mientras que de 35 grados en adelante se estabilizan cerca de 1.36. A esto se suman el grosor del rayo, que en la escala ocupaba casi un grado, y la dificultad de leer a oscuras.

La serie de la cara curva fue más confiable por una razón sencilla: ahí el ángulo de salida crece muy rápido, de modo que un error de un grado en la lectura pesa mucho menos.

El ángulo crítico también salió bien. Con el índice de la cara curva, la ecuación \eqref{eq:critico} da
\begin{equation}
\theta_c = \arcsin\left( \frac{1}{1.494} \right) = 42.0 \text{ grados},
\label{eq:critico_calc}
\end{equation}
y en el experimento el rayo transmitido todavía existía a 40 grados, casi rasante a 74 grados, y ya no a 45. El valor medido queda entre esos dos ángulos, justo donde debe.

Hay un detalle curioso más. En la serie de la cara plana, al llegar la luz casi rasante a 89 grados, el rayo dentro de la lente marcó 48 grados. Ese ángulo es el máximo que puede tener un rayo refractado y debería coincidir con el ángulo crítico de 42 grados. La diferencia de 6 grados apunta otra vez al mismo problema de alineación.

En conclusión, la relación $n_i \sin\theta_i = n_t \sin\theta_t$ describió muy bien el comportamiento de la luz en ambos sentidos, con ajustes lineales de $R^2 = 0.9972$ y $R^2 = 0.9997$. El mejor valor del índice de refracción del acrílico fue 1.494, con apenas 0.2 por ciento de error, y se observó la reflexión total interna a partir de un ángulo crítico cercano a 42 grados, en el que además el rayo reflejado cumplió la ley de la reflexión.

La práctica también dejó una lección. Una buena alineación importa tanto como una buena lectura.
